\documentclass[10pt]{article}
% tex/env.tex — Packages and environment configuration
%
% Package and environment choices are aligned with the gold-standard reference
% material under `tex/literature/` (e.g., TRM and trust-region analyses).

\usepackage[utf8]{inputenc}
\usepackage[T1]{fontenc}
\usepackage{lmodern} % scalable Latin Modern fonts (avoids bitmap EC fonts)

% Typography and figures/tables (reference-style baseline)
\usepackage{microtype}
\usepackage{graphicx}
\usepackage{subcaption}
\usepackage{booktabs}
\usepackage{multirow}

% Math and theorem tooling
\usepackage{amsmath,amssymb,amsthm}
\usepackage{mathtools}
% `thmtools` is not needed by the paper body.  TeX Live 2026's thmtools
% patch currently conflicts with amsthm's shared-counter declarations, so we
% rely directly on amsthm for the theorem environments below.
\allowdisplaybreaks

% Algorithms
\usepackage{algorithm}
\usepackage{algorithmic}

% Colors and links/cross-references
\usepackage[dvipsnames]{xcolor}
\usepackage[hyperfootnotes=false]{hyperref}
\usepackage[capitalize,noabbrev]{cleveref}

% Appendix references:
% We label appendix entries via \label[appsec]{app:...} so that \Cref{app:...}
% prints "Appendix ..." rather than "Section ...", matching common paper style.
\crefname{appsec}{Appendix}{Appendices}
\Crefname{appsec}{Appendix}{Appendices}
% cleveref does not infer names for custom theorem environments.
\crefname{assumption}{Assumption}{Assumptions}
\Crefname{assumption}{Assumption}{Assumptions}

% Page and bibliography
\usepackage[margin=1in]{geometry}
\usepackage{natbib}

% Lists and code listings
\usepackage{enumitem}
\usepackage{listings}

% Attempt to make hyperref and algorithm work together better (reference convention).
% Recent TeX Live releases define this hook in `algorithm` already, so provide
% the fallback only when the package has not defined it.
\providecommand{\theHalgorithm}{\arabic{algorithm}}

% Development-only optional packages (guarded to avoid build failures).
\IfFileExists{todonotes.sty}{\usepackage[textsize=tiny]{todonotes}}{}
% tcolorbox is optional and is not used by the paper body.  Some TeX Live
% installations ship tcolorbox without the auxiliary tikzfill.image module
% required by the `most` library; guard both files so the document remains
% portable in such environments.
\IfFileExists{tcolorbox.sty}{%
  \IfFileExists{tikzfill.image.sty}{\usepackage[most]{tcolorbox}}{}%
}{}

%==============================================================================
% THEOREM ENVIRONMENTS (shared counter, grouped by style)
%------------------------------------------------------------------------------
% The structure matches the reference material conventions (e.g., TRM), and is
% intentionally kept in `env.tex` so all sections share the same definitions.
%==============================================================================
\theoremstyle{plain}
\newtheorem{theorem}{Theorem}[section]
\newtheorem{lemma}[theorem]{Lemma}
\newtheorem{proposition}[theorem]{Proposition}
\newtheorem{corollary}[theorem]{Corollary}

\theoremstyle{definition}
\newtheorem{definition}[theorem]{Definition}
\newtheorem{assumption}[theorem]{Assumption}
\newtheorem{example}[theorem]{Example}

\theoremstyle{remark}
\newtheorem{remark}[theorem]{Remark}

% Listing style for Python code
\lstset{
  language=Python,
  basicstyle=\ttfamily\small,
  keywordstyle=\color{blue},
  commentstyle=\color{gray},
  frame=single,
  breaklines=true,
  breakatwhitespace=true,
  xleftmargin=0pt,
  escapeinside={<@}{@>}
}

\input{math}
\geometry{margin=0.72in}
\title{REINFORCE Pro Max: Causal Prefix Filtering and
Adaptive Advantage Normalization for Critic-Free LLM-RL}
\author{Anonymous}
\date{}

\begin{document}
\maketitle

\begin{abstract}
Critic-free reinforcement learning for language models must control both
token-mass imbalance in sampled advantages and the mismatch between rollout and
training policies. We present REINFORCE Pro Max, which combines leave-one-out
advantages, sign-preserving asymmetric normalization, and a prefix-cumulative
importance-sampling gate. The gate is computed from the cached rollout/snapshot
ratio but gradients flow only through the current/snapshot PPO ratio. On a
finite autoregressive tree, we prove an exact decomposition of the implemented
standard-PPO objective into the raw surrogate, gate and normalization
corrections, and a nonnegative clipping deduction. A finite-tree adaptive
bound then gives a conditional sufficient certificate for reward improvement.
The mathematical statements and their finite-tree proofs are checked in Lean;
the formalization does not claim to verify floating-point neural-network
training. The current artifact contains structural synthetic checks only and
does not report an LLM benchmark.
\end{abstract}

\section{Problem and scope}
Critic-free objectives such as RLOO and GRPO are widely used for reasoning
tasks \citep{ahmadian2024back,shao2024deepseekmath}, while PPO remains a
standard clipped policy-gradient baseline \citep{schulman2017proximal}. Recent
long-horizon trust-region analyses motivate retaining per-position causal
structure \citep{li2025trm}.
LLM-RL rollouts are generated by $\pi_{\rm roll}$ while the actor is updated
from a frozen snapshot $\pi_{\rm old}$. For response $y=(y_1,\ldots,y_T)$,
write $\rho_t=\pi_\theta(y_t\mid c_t)/\pi_{\rm roll}(y_t\mid c_t)$ and
$w_t=\pi_{\rm old}(y_t\mid c_t)/\pi_{\rm roll}(y_t\mid c_t)$. The analyzed
configuration is standard PPO (\texttt{policy\_loss\_type=ppo}),
\texttt{vllm\_is\_correction\_type=reinforce\_pro}, no dual clip, no KL or entropy
term, and sparse terminal reward with zero KL coefficient. These restrictions
are essential: optional dual clipping and auxiliary losses change the objective.

We study a finite response alphabet and finite horizon with mutual support,
bounded reward $|R|\le R_{\max}$, and explicitly stated trust-region or
moment assumptions. The results are conditional certificates, not guarantees
that an arbitrary trained neural policy satisfies those assumptions.

\section{Method}
For a prompt group of $n$ responses, the raw leave-one-out signal is
$A_i=R_i-(n-1)^{-1}\sum_{j\ne i}R_j$. After token expansion, positive and
negative active tokens need not have equal mass. Let $P,N$ be the positive and
negative token sets. The ideal ProMax factors are chosen so that
$\sum_{P}\alpha A+\sum_N\beta A=0$ and
$|P|^{-1}|N|^{-1}$-free second-moment normalization equals one on nonzero
tokens, with $\alpha,\beta>0$. Thus signs and within-sign ordering are
preserved. The implementation adds an epsilon, cap, clamp, and fallback; the
formal result treats their moment effects as explicit residuals.

For each active token, the prefix proxy is
\[
 M_t^{\rm pre}=\mathbf 1\!\left\{
 \ell_-\le \frac{1}{t}\sum_{s\le t}\log w_s\le \ell_+\right\},
 \qquad w_t=\exp(\log\pi_{\rm old}-\log\pi_{\rm roll}).
\]
The implemented loss is
\[
 \mathcal L=\operatorname{MaskedMean}_t\left[
 M_t^{\rm pre}\,\operatorname{stopgrad}(w_t)
 \left(-\min(r_t\hat A_t,\operatorname{clip}(r_t)\hat A_t)\right)\right],
\quad r_t=\frac{\pi_\theta}{\pi_{\rm old}}.
\]
The prefix gate, cached importance coefficient, action mask, and normalized
advantage are detached. Only $r_t$ carries policy gradient.

\section{Formal results}
Let $k$ index represented response/token positions, let $\omega_k\ge0$ be the
chosen reduction weight, $M_k$ the gate, $r_k$ and $w_k$ the two probability
ratios, $A_k$ the raw signal, and $H_k$ the implemented normalized signal. For
standard PPO define
\[
 U=\sum_k\omega_k r_kw_kA_k,\quad
 G=\sum_k\omega_k r_kw_k(1-M_k)A_k,\quad
 N=\sum_k\omega_k r_kw_kM_k(H_k-A_k),
\]
and
\[
 C=\sum_k\omega_kM_kw_k\bigl(r_kH_k-
 \min(r_kH_k,\operatorname{clip}(r_k)H_k)\bigr).
\]
Then the finite-batch identity is
\[
 -\mathcal L=U-G+N-C,\qquad C\ge0.
\]
The gate and normalization terms are signed corrections. The nonnegativity of
$C$ is a standard-PPO fact and does not extend to dual clip; a Lean-checked
counterexample is included in \texttt{formal/ProMaxObjective.lean}.

For a finite autoregressive tree, if every rollout-reachable history of length
less than $T$ satisfies either a local TV bound $\varepsilon$ or a local
reverse-KL bound $\delta$, the future continuation discrepancy after history
length $t$ is bounded by
\[
 \min\{1,(T-t)\varepsilon,\sqrt{(T-t)\delta/2}\}.
\]
With bounded rewards, summing these continuation errors bounds the difference
between the raw surrogate and the true reward change. Consequently, a
corrected empirical ProMax gain exceeding the explicit gate, normalization,
clipping, sampling, and adaptive remainder terms is sufficient for positive
true reward gain. The statement is for independent evaluation blocks and a
fixed finite candidate family, or a pre-specified class with a fixed finite
cover.

The finite-tree identities, adaptive recursion, RLOO moments, reduction
bounds, and binary-length covariance conditions are checked by the Lean files
listed in \texttt{formal/PROOF\_AUDIT.md}. The audit currently covers 31 labeled
mathematical environments and 911 project theorems.

For binary rewards, the ideal normalization has two class amplitudes whose
ratio is the ratio of class-conditional mean response lengths. With the actual
epsilon stabilizer, cap, clamp, fallback, and token denominator, the complete
on-policy update remains a positive multiple of the success-probability
gradient when active length depends only on the reward class and both classes
have positive mass. If length also varies within a reward class, the update
decomposes into this class multiplier plus a within-class covariance error; a
strict local-improvement conclusion then requires an explicit covariance
margin. The theorem does not claim that arbitrary content-dependent lengths or
an arbitrary Adam step improve reward.

\section{Implementation correspondence}
The code path uses active-token masking and distinguishes the differentiable
PPO ratio from the detached rollout correction. Standard PPO is fixed to the
configuration above; GSPO, dual clip, KL loss, entropy loss, MoE auxiliary
loss, reward shaping, and dynamic group filtering require separate objective
terms and are outside the stated certificate. Variable-length dynamic
microbatch response-count weighting does not automatically equal a global
token mean, so the population reduction must be reported explicitly.

\section{Evaluation status and limitations}
The current repository includes deterministic synthetic structural checks for
normalization and prefix masking. It does not contain matched LLM training,
multiple seeds, an independent AIME/MATH-style test set, or a benchmark claim.
The formal proofs use exact real arithmetic on finite models; they do not
formalize floating-point overflow, GPU kernels, Adam, Ray/vLLM execution, or
the assumptions of a general neural policy. These are required before claiming
empirical ICLR readiness.

\bibliographystyle{plainnat}
\bibliography{main/reference}
\end{document}
